\begin{onehalfspace}

Questa appendice documenta il classificatore valutato nella prima fase
sperimentale (Capitolo~\ref{chap:risultati}) e, soprattutto, separa due
domande che nella prima versione della campagna erano state confuse:
\emph{il candidato più frequente conduce ai fattori?} (TOP-1) e
\emph{esiste un candidato utile fra i primi sedici?} (TOP-16). Tutti i
numeri riportati provengono dagli artefatti schema~2 rigenerati con il
circuito $N=15$, $a=7$ corretto, il medesimo manifest del circuito e seed
deterministici.

\subsection{L'Ipotesi alla Prova: il Classificatore ML}
\label{sec:risultati_uc_m2}
\label{sec:valutazione_clf}

Il Metodo~2 usa un classificatore come \emph{gate di qualità}: riceve la
distribuzione normalizzata sui $2^8=256$ esiti e, se approva
l'istogramma, tenta l'estrazione dei fattori dai quattro candidati più
frequenti. L'ablazione $M_{\mathrm{TOP4}}$ applica la stessa ricerca senza
filtro e consente quindi di distinguere l'effetto del classificatore da
quello della ricerca multi-candidato.

Per ciascuno dei due scenari sono stati generati $2\,000$ istogrammi da
$1\,024$ shot. I cinque parametri del modello uniforme
($\lambda_{1q}$, $\lambda_{2q}$, $T_1$, $T_2$ e readout) sono stati
campionati indipendentemente in un intervallo relativo del $\pm50\%$
attorno al preset. Il dataset è stato separato in train, validation e
test con proporzioni $60/20/20$: la scelta dell'architettura usa soltanto
la validation; il test viene valutato una sola volta dopo il refit su
train+validation.

La regola TOP-1 produce $1\,262$ positivi e $738$ negativi in UC1
($63.1\%/36.9\%$), e $1\,453$ positivi e $547$ negativi in UC2
($72.7\%/27.3\%$). Random Forest, \ac{SVM} RBF calibrata e \ac{MLP}
$(256,128)$ sono state confrontate mediante F1 sulla validation. La
Tabella~\ref{tab:clf_metrics} riporta la valutazione indipendente del
modello selezionato.

\begin{table}[H]
\centering
\caption{Metriche TOP-1 sul test set indipendente (400 istogrammi per scenario)}
\label{tab:clf_metrics}
\begin{tabular}{llccc}
\toprule
\textbf{Scenario} & \textbf{Modello selezionato} & \textbf{F1} &
\textbf{Accuratezza} & \textbf{\acs{AUC}-ROC} \\
\midrule
UC1 & \ac{SVM} & 0.919 & 0.895 & 0.965 \\
UC2 & \ac{SVM} & 0.923 & 0.888 & 0.958 \\
\bottomrule
\end{tabular}
\end{table}

Sulla validation l'\ac{SVM} ottiene F1 pari a $0.908$ in UC1 e $0.929$
in UC2, contro $0.811/0.842$ della Random Forest e $0.874/0.906$
dell'\ac{MLP}. Queste metriche mostrano che il modello riproduce bene la
regola TOP-1; non dimostrano però che il filtro riduca il costo operativo.
Quel confronto richiede gli stessi istogrammi per M1, TOP-4 e M2 ed è
riportato nella Sezione~\ref{sec:confronto}.

\begin{figure}[H]
    \centering
    \includegraphics[width=0.85\textwidth]{fig_roc_curve.pdf}
    \caption[Curve ROC dei classificatori TOP-1 selezionati]{Curve ROC
    sul test set indipendente dei due classificatori \acs{SVM} selezionati
    mediante validation. Le aree sono $0.965$ (UC1) e $0.958$ (UC2); la
    diagonale rappresenta un classificatore casuale.}
    \label{fig:roc_curve}
\end{figure}

% -----------------------------------------------
\subsection{Audit delle Etichette: TOP-1 e TOP-16 Sono Problemi Diversi}
\label{sec:audit_etichette}

La campagna rigenerata calcola entrambe le etichette sullo \emph{stesso}
istogramma, eliminando l'ambiguità del dataset storico. L'esito è
riassunto nella Tabella~\ref{tab:audit_etichette}.

\begin{table}[H]
\centering
\caption[Audit delle etichette TOP-1 e TOP-16]{Bilanciamento delle due
regole sul dataset completo. TOP-16 è a classe unica e non consente un
addestramento discriminante.}
\label{tab:audit_etichette}
\begin{tabular}{lrr}
\toprule
\textbf{Regola} & \textbf{UC1: positivi/negativi} &
\textbf{UC2: positivi/negativi} \\
\midrule
TOP-1  & $1\,262/738$ & $1\,453/547$ \\
TOP-16 & $2\,000/0$   & $2\,000/0$ \\
\bottomrule
\end{tabular}
\end{table}

Il runner non forza un modello sul dataset TOP-16: registra
esplicitamente \texttt{single-class-dataset} e interrompe
l'addestramento per quella regola. M2 usa perciò soltanto i modelli TOP-1
e conserva nel payload \texttt{label\_top\_k=1}. Questo controllo evita
di attribuire capacità predittiva a un classificatore costante.

Una diagnostica separata su 60 esecuzioni nominali UC1 chiarisce la
struttura della classe TOP-1. La moda è $y=0$ in $30/60$ casi e conduce
ai fattori negli altri $30/60$; i primi quattro candidati contengono
invece almeno un esito utile in $60/60$ casi. La Tabella
\ref{tab:struttura_negativi} riporta la distribuzione della moda.

\begin{table}[H]
\centering
\caption[Distribuzione della moda su UC1]{Distribuzione della misura più
frequente su 60 istogrammi nominali indipendenti.}
\label{tab:struttura_negativi}
\begin{tabular}{lrr}
\toprule
\textbf{Valore della moda} & \textbf{Occorrenze} & \textbf{Frazione} \\
\midrule
$0$   & 30 & $50.0\%$ \\
$64$  &  5 & $8.3\%$ \\
$128$ & 15 & $25.0\%$ \\
$192$ & 10 & $16.7\%$ \\
\bottomrule
\end{tabular}
\end{table}

L'esito $y=0$ rappresenta la fase nulla e non fornisce il periodo, ma
non implica che l'istogramma sia privo di segnale. Nel campione
rappresentativo della Figura~\ref{fig:audit_istogramma_negativo}, la
moda $y=0$ raccoglie 176 conteggi, mentre $y=64$, $192$ e $128$ ne
raccolgono rispettivamente 160, 141 e 131: tutti e tre conducono ai
fattori. La massa complessiva sui quattro picchi è $0.594$.

\begin{figure}[H]
    \centering
    \includegraphics[width=0.86\textwidth]{audit_istogramma_negativo.pdf}
    \caption[Istogramma TOP-1 negativo ma recuperabile con TOP-4]{Esempio
    nominale UC1 in cui TOP-1 fallisce perché la moda è $y=0$, mentre i
    successivi tre picchi conducono ai fattori. Il classificatore TOP-1
    apprende quindi una regola sul candidato dominante, non una misura
    generale della presenza di segnale quantistico.}
    \label{fig:audit_istogramma_negativo}
\end{figure}

\subsubsection{Perché TOP-16 Produce una Classe Quasi Sempre Positiva}
\label{subsec:soglia_classe_negativa}

Per $N=15$, $a=7$ e otto bit di controllo, 63 dei 256 esiti conducono ai
fattori tramite la procedura di frazioni continue adottata. Anche per un
istogramma uniforme, la probabilità che nessuno dei primi $K$ candidati
sia utile è quindi
\begin{equation}
    P_{\mathrm{neg}}(K)=
    \frac{\binom{256-63}{K}}{\binom{256}{K}}.
    \label{eq:tetto_negativi}
\end{equation}
Per $K=16$ essa vale soltanto $0.9275\%$: la regola assegna dunque
un'etichetta positiva a oltre il $99\%$ degli istogrammi uniformi per
puro effetto combinatorio. Il risultato $2\,000/0$ osservato nei due
dataset è coerente con questo limite e non autorizza a concludere che il
segnale sia perfetto.

\begin{table}[H]
\centering
\caption[Tetto combinatorio alla classe negativa in funzione di $K$]{Probabilità
di non includere alcuno dei 63 esiti utili scegliendo $K$ celle da 256
senza informazione sulla distribuzione.}
\label{tab:tetto_topk}
\label{tab:soglia_negativi}
\begin{tabular}{cc}
\toprule
\textbf{Regola} & $P_{\mathrm{neg}}(K)$ \\
\midrule
TOP-1  & $75.39\%$ \\
TOP-2  & $56.76\%$ \\
TOP-4  & $32.06\%$ \\
TOP-16 & $0.93\%$  \\
\bottomrule
\end{tabular}
\end{table}

\begin{figure}[H]
    \centering
    \includegraphics[width=0.78\textwidth]{tetto_classe_negativa.pdf}
    \caption[Tetto combinatorio alla classe negativa]{Decadimento di
    $P_{\mathrm{neg}}(K)$ al crescere del numero di candidati. La figura
    è rigenerata dai dati schema~2 e rende visibile perché TOP-16 non
    definisce un problema di classificazione bilanciato su questa
    istanza.}
    \label{fig:tetto_classe_negativa}
\end{figure}

Il verdetto ha un perimetro preciso: non mostra che l'apprendimento
automatico sia inutile in generale, ma che un filtro a valle addestrato
su TOP-1 replica una decisione già superata da TOP-4, mentre la regola
TOP-16 è quasi sempre positiva per costruzione. Questo motiva lo
spostamento dell'apprendimento dalla distribuzione finale alla sindrome
di errore nei capitoli dedicati alla \ac{QEC}.

\end{onehalfspace}
